\section{Введение: как устроен курс}

\subsection*{О чём эта книга}

Эта книга учит программировать квадрокоптер «Геоскан Пионер»: от первой строчки на языке Lua до автономной миссии, в которой дрон сам взлетает, облетает заданные точки, реагирует на нештатные ситуации и садится. Вторая половина книги делает то же самое на языке Python, а затем выходит за пределы симулятора --- к камере, визуальным маркерам и групповым полётам на настоящем дроне.

Главный принцип курса --- \textbf{каждый пример запускается}. Все программы из книги лежат в папке \texttt{code/} рядом с текстом и автоматически проверяются в симуляторе: взлетают, пролетают свои точки и садятся. Если пример в книге не совпадает с тем, что вы видите на экране, --- это ошибка книги, а не ваша.

\subsection*{Где писать код}

Для занятий нужен только браузер: веб-симулятор \url{https://simulator.sakoryagin.ru} показывает полёт в 3D, умеет Lua, Python и Blockly и использует тот же API, что настоящий Пионер. Настоящий дрон понадобится ближе к концу курса --- для камеры и компьютерного зрения. Как запускать программы в симуляторе и на дроне, объясняет глава \ref{sec:start}.

\subsection*{Как читать главы}

Каждая глава построена одинаково:
\begin{enumerate}
    \item \textbf{Цели} --- что вы будете уметь после главы.
    \item \textbf{Теория с примерами} --- короткие программы, которые нужно запустить и сверить результат с комментариями в коде.
    \item \textbf{Проверьте себя} --- вопросы с ответами.
    \item \textbf{Практика} --- задания без готового решения; решения к ним лежат в \texttt{code/lua/tasks/}.
    \item \textbf{Итоги} --- что запомнить.
\end{enumerate}

В тексте встречаются рамки:
\begin{itemize}
    \item \textcolor{GeoskanBlue}{\faInfoCircle}~\textbf{Информация} --- пояснение или ссылка на другой раздел;
    \item \textcolor{red!70!black}{\faExclamationTriangle}~\textbf{Внимание} --- типичная ошибка или вопрос безопасности;
    \item \textcolor{GeoskanOrange}{\faLightbulb}~\textbf{Подсказка} --- приём, который экономит время;
    \item \textcolor{GeoskanDarkBlue}{\faGavel}~\textbf{Правило} --- соглашение, которого курс придерживается всегда.
\end{itemize}

\subsection*{Дорожная карта}

\begin{enumerate}
    \item \textbf{Подготовка.} Интерфейс симулятора, первая программа, настоящий Пионер и техника безопасности.
    \item \textbf{Часть I. Lua.} Язык (переменные, условия, циклы, функции) $\to$ светодиоды и таймеры $\to$ события автопилота и автономные миссии.
    \item \textbf{Часть II. Python.} Основы языка $\to$ библиотека \texttt{pioneer\_sdk}: подключение, светодиоды, взлёт, полёт по точкам $\to$ камера, маркеры ArUco, трекинг, рой.
    \item \textbf{Часть III. Профессиональная навигация.} Теория визуальных маркеров, калибровка камеры, карты, фильтрация и точная посадка.
    \item \textbf{Приложения.} Справочник Lua API, шпаргалка Lua~$\leftrightarrow$~Python, настройка пульта, протокол безопасности.
\end{enumerate}

Lua и Python в курсе идут в одном порядке: сначала язык, потом светодиоды, потом полёт. Поэтому главу про Python можно читать и отдельно, но полезнее --- после Lua: событийная модель автопилота, разобранная на Lua, объясняет, почему Python-программа «ждёт точку» именно так.

\subsection*{Планируемые результаты}

После курса вы сможете:
\begin{itemize}
    \item писать читаемый код на Lua и Python и объяснять, что делает каждая строка;
    \item управлять бортовой периферией --- светодиодами и таймерами --- без блокировки программы;
    \item строить программу полёта как конечный автомат: команда $\to$ событие $\to$ следующая команда;
    \item предусматривать аварийные ситуации: удар, разряд батареи, потеря связи;
    \item переносить программу из симулятора на настоящий дрон без переписывания.
\end{itemize}
